\mychapter{Relevância do Trabalho}{cap:relevancia}
\lhead{RELEVÂNCIA DO TRABALHO}

A contribuição principal recai sobre a computação, e não sobre o domínio das aplicações de exemplo. O trabalho leva para o desktop a pergunta que SequalsK e GUIMigrator respondem apenas para Android e iOS e busca endereçar uma limitação que o próprio SequalsK aponta: a tradução da camada de interface, que limitou o código compartilhado a 16\% em um dos seus estudos de caso \cite{schultes2025native}. Na literatura triada, há uma lacuna de contexto (poucas avaliações em desktop com interface nativa) e metodológica, pois as comparações de desempenho e os estudos de tendências encontrados tratam de frameworks móveis \cite{biorn2020empirical, jost2025state}.

O trabalho entrega evidência empírica sobre três abordagens avaliadas nas mesmas aplicações de referência: código nativo escrito à mão, geração a partir de fonte única e abordagem multiplataforma convencional. Para desenvolvedores que valorizam desempenho e aderência ao design de cada sistema, os resultados ajudam a decidir entre manter uma base de código por sistema, com custo alto, e adotar um framework multiplataforma, com interface que destoa da plataforma. Na abordagem proposta, por gerar projetos nativos padrão, o código permanece inspecionável em cada sistema, o que tende a reduzir o aprisionamento a uma camada opaca.

O trabalho também adota geração determinística, em contraste com a migração por modelos de linguagem. Tanto o GUIMigrator quanto o conversor de Mahmoud et al.\ relatam resultados superiores aos de \textit{baselines} baseados em LLMs \cite{gao2026guimigrator, mahmoud2024trans}. A geração de interfaces a partir de modelos já é uma linha estabelecida \cite{ruiz2019evaluating}, e aqui ela é aplicada à produção de componentes nativos de desktop.

No contexto de desktop, estudos comparativos como o de Milichiewicz e Badurowicz avaliam ferramentas multiplataforma (Electron, Tauri e PyQt) entre si e apontam compromissos de integração com o sistema e consumo de recursos \cite{milichiewicz2025comparative}, mas não incluem uma versão nativa desenvolvida manualmente como \textit{baseline}. A novidade, portanto, está na combinação de geração a partir de fonte única, interface nativa de desktop e comparação com o código nativo.

O trabalho é viável no prazo, pois o escopo se restringe a aplicações simples para Windows e macOS, há equipamentos disponíveis e o projeto das aplicações dispensa dados externos.